% Section 9. Grounded rewrite 2026-11-03 (branch docs/cep-grounded-draft). Only numbers stated in sections 2-8, with
% their scripts there. Claim changes: papers/cep/CLAIM_CHANGES.md.
\section{Discussion and limitations}\label{sec:discussion}

\subsection{What the installed camera delivers now}\label{sec:performs}

The optical side is settled up to its stated inputs. The scale at the focused plane is known per axis to about
$\pm 10$\,\% at 95\,\% (\tabref{tab:scale}), with a discrete $+8.9$\,\% open if the stated tube length includes a
CS-to-C ring. The plant's ``$\approx 6$~cm'' for the tube dominates the budget; every other term is about 2\,\% or
less. Any detector on the clips must use the anisotropic clip pixel, and its resolution element is the PAL sample in
both recordings (\secref{sec:budget}).

At F/1.6 the camera samples a slice of $\pm 9.4$~\um{} about the focused plane (\tabref{tab:dof}). Within it the
magnification is constant well below the scale uncertainty (mean depth correction 0.005--0.013\,\%;
\secref{sec:depthsmall}), so a single focal-plane scale suffices. The same thinness makes the size prior
uninformative (\secref{sec:disamb-result}). Because the size ratio depends only on $\delta / d_\mathrm{focus}$
(Eq.~\ref{eq:ratio}) while the admitted depth range grows in proportion to $N$ (Eq.~\ref{eq:coc}), stopping down to
F/16 widens the slice tenfold to $\pm 94$~\um{} and still leaves the admitted range below the
$\delta \approx 1.3$~mm at which the ratio reaches 5\,\%.

Neither classical detector measures crystal size on the case-study clip
(\S\S\,\ref{sec:csd-edge}--\ref{sec:watershed}). Resolution is not the obstacle for most of a strike: from about half
an hour after seeding the microscope medians exceed 0.1~mm (\tabref{tab:microscope}), tens of PAL samples, and only
the slurry crystals of 10--19~\um{} sit at the resolution limit (\secref{sec:expected}). The missing piece is a
detector validated against hand-labelled crystals (\secref{sec:validation}).

\subsection{Limitations}\label{sec:notperform}

The scale is modelled, not measured in the image: no object of known size has been imaged through the glass. The
tube length, air gap, glass thickness and F-number are approximate plant statements, the camera model with its cell
size and CCD clock is the best-matching candidate, to be confirmed, and the 13.5~MHz sampling is the BT.601 rate, not
a datasheet value of the installed recorder (\tabref{tab:inputs}). The case-study clip has no size reference: its
recording date is contradicted by the image data, so its strike phase is unknown (\secref{sec:recording}), and the
microscope series is from another day. That series consists of chords of operator-picked crystals, read once by an AI
transcription from the photographs' overlays, not a measured distribution (\secref{sec:expected}). The depth of field
rests on two choices, the resolution element (\tabref{tab:dof}) and the PSF constant $k = 1/2$, which has no
page-checked source and sets the slice but not the asymptotic exponent (\secref{sec:psf}). The sight glass is treated
paraxially, its residual distortion over the field has not been computed, and its index is assumed
(\secref{sec:shift}). Slurry homogeneity is untested (\secref{sec:selfcheck}). Frame content matters: the recorder
overlays and the pillarbox move the sharpness slope from 5.36 to 4.88--4.90 (\tabref{tab:crop}), and the right picture
edge moves by about 8 clip pixels between and within clips, so crops must keep a margin (\secref{sec:recording}).

\subsection{Operational recommendations}\label{sec:recs}

\begin{enumerate}
\item \textbf{Measure the geometry at the next shutdown}: read the camera's model sticker (cell size, CCD clock),
caliper the extension tube or spacer and photograph the mount (CS-to-C ring), and read the recorder model (sampling)
and the glass stamp (index). A ruler or graticule imaged at the inner glass face of the empty pan would measure the
scale in the image and locate the focal plane.
\item \textbf{Record the PAL stream, not the monitor.} Screen recordings add no resolution but an anisotropic
stretch and a moving picture edge. Logging the recorder clock against the control system keeps the time since
seeding traceable; on the 2026-02-04 recording the two agree to about $\pm 30$~s (\code{scripts/check_anchors.py}).
\item \textbf{Replace the classical detectors by a validated segmentation}, such as a fine-tuned Cellpose model
scored as in \secref{sec:validation}, with a size window that admits the largest microscope chords, which exceed the
current 1~mm bound (\secref{sec:expected}).
\item \textbf{Treat the camera as a thin-slice sensor at F/1.6.} A single focal-plane scale suffices; depth
disambiguation would need detections far outside the in-focus slice, which neither the engraved stops nor the current
detectors provide, and stopping down trades the slice against light and diffraction (\secref{sec:dof}).
\end{enumerate}

\subsection{Towards a soft sensor}\label{sec:softsensor}

The control use of the camera is a soft sensor of the growth state of each strike, with the growth model of
\secref{sec:growth} as measurement equation, anchored in physical units by the microscope series. Driven by a
validated camera median against time since seeding, a state estimator could return the time since seeding or flag a
strike that leaves the nominal curve, and be compared with an estimator using the plant's process signals alone. This
paper provides the scale, the depth model and the validation protocol; the camera size signal is
\pending{validated camera size against time, 2026-10-30}.
